% language=us runpath=texruns:manuals/moretocome

\startcomponent moretocome-optimizations

\environment moretocome-style

\startchapter[title=Optimizations]

I've written about performance before, and occasionally I come back to it. We're
now in 2025 and we can safely say that performance of the \LUAMETATEX\ plus
\CONTEXT\ is quite okay. Users who also use other macro packages often express
how much faster \CONTEXT\ is but I can't really check that: I don't have other
macro packages installed, it would also depend on the installations and of course
the kind of documents matter too. However, when it comes ot \LUATEX\ and
\LUAMETATEX, I'm not that surprised if \CONTEXT\ is actually relatively fast;
after all these engines originate in the \CONTEXT\ community.

There are, to summarize several factors that influence performance. Start by
asking yourself \quotation {What is left to \TEX\ and what is done in \LUA?} and
\quotation {How good are the \TEX\ macros (and solutions) that I use and how
efficient is the \LUA\ code, if used?}. These are subjective matters so best left
aside. I just like to remark that when you read comments on the engines or
\CONTEXT\ by those not too familiar with it, you should take them with a grain of
salt: one can read plenty of non-sense.

We can discuss where in the engine or \CONTEXT\ we try to make sure that
performance is good, but we can also start from the other end. What are actually
the potential bottlenecks. Some are just side effects of design decisions, the
nature of the problem, or the way things are implemented, for instance expansion
extensive macros solutions. There's also user interfacing to take into account:
we don't want to compromise there. When we started with \TEX\ and reached a state
where we could qualify \CONTEXT\ as a macro package, sometime mid 90's, for sure
the user interface came with a performance penalty.

When we moved on to \LUATEX\ in the first decade of this century, the low level
implementation of the user interface was rewritten. On the one hand the memory
footprint got smaller, inheritance more clever, but in theory at a larger
expansion cost. However in practice performance got a bit better. Then, in the
decade when we went for \LUAMETATEX\ extensions to the macro (preamble) parsing
capabilities of the engine made for a better performance again. If you realize
what goes on in the user interface one can actually be surprised that it is so
fast at all but course computers also got faster.

So why do I try (or even look at) potential optimizations every now and then? A
valid reason is an observed slow-down, or a future slow-down due to evolving
features. But \quotation {Wait!} you say, \quotation {Isn't \CONTEXT\ in general
becoming faster over time?}. Indeed that is the case, but a side effect of adding
more advanced feature is that the can actually add a little runtime. We're
talking percentages but they can accumulate. So, when I am in for a little
gaming, I look at solutions, at features that were added but never fully applied,
and surprise: we sometimes can gain back a few percent. Of course there are
exceptions, for instance when you enable tagging, which recently has become a bit
more popular, runtime can increase (sometimes double) but those are functionality
hound so there is little to gain: you're asking for something that cost time. But
even there we might gain back a little over time.

One problem with noticing a gain in performance is that a move to a different
machine or upgrade of the operating system can obscure observations. Another
interference comes from additional features and enabling them. So in practice
differences in performance can only be noticed when a document is relatively
stable (a few pages more hardly change runtime) and a similar version or style is
used. It is also a bit feeling: did it get slower or faster, a perception.

So, how important is going from 15.2 seconds down to 14.7 on the \LUAMETATEX\
manual really, or 7.1 down to 6.9 in our math reference book. I have a laptop
with a 2017 (Xeon) processor, and a 2025 with that years fastest AMD processor
could more than half that time. So why not just buy a new laptop? Valid reasons
are that \TEX\ doesn't pay the bills so there is no pay-back, why waste a
properly working laptop and put a burden on the environment, and installing a new
machine is no fun either. So why not just optimize a bit and stick to the machine
as long as possible? A side effect of a little optimization here and there, based
on observed usage patterns also has the side effect of (maybe) smoother runs on
virtual machines that compete with each other.

End 2025 I decided to look a bit at the potential bottlenecks of some documents
again saw confirmed that attributes are part of the drag. Attributes are one of
the first extensions to \LUATEX. When set, they travel with a node. For instance,
when a character in the input becomes a glyph node, it gets a list of currently
set attributes.

\starttyping
[glyph 'a'] <=> [glyph 'b'] <=> [glyph 'c']

[attr 1]        [attr 1]        [attr 1]
[attr 3]        [attr 3]        [attr 2]
[attr 8]        [attr 8]        [attr 6]
\stoptyping

These attributes are kept in a list, like:

\starttyping
[attr 1] -> [attr 3] -> [attr 8]
[attr 1] -> [attr 3] -> [attr 8]
[attr 1] -> [attr 2] -> [attr 6]
\stoptyping

And because it is likely that lists the same for successive nodes, we actually
have:

\starttyping
[attr list 1] -> [attr 1] -> [attr 3] -> [attr 8]
[attr list 1] -> [attr 1] -> [attr 3] -> [attr 8]
[attr list 2] -> [attr 1] -> [attr 2] -> [attr 6]
\stoptyping

where the first two are the same list. Every list has a reference counter so here
we have two lists, one with a reference count of two, and one with a count of
one. These attributes are small nodes themselves but in \CONTEXT\ on a page we
have thousands of them, and in for instance balanced pages, as in column sets, we
have tens of thousands. We need to allocate them and at some point they might get
tested for every node in a list. In fact, if these attribute lists have 10
elements, and we have 4000 nodes on page, checking for an attribute checking 4000
attributes lists with in a worst case scenario the wanted attribute at the end.

There are various optimizations. For instance we could avoid a check by checking
if the attribute list itself changed (when we work on neighboring nodes). The
lists are sorted so a lookup is optimized by quitting a lookup when we passed the
asked attribute index. List construction is also optimized, which in itself is
kind or tricky if you take grouping into account. here \ \LUAMETATEX\ is more
advanced than \LUATEX, which could be one of the reasons that it is faster than
\LUATEX.

A recent (2025) optimization in lookups is that at the cost of extra
housekeeping, we don't need to check if an attribute is not in a list. A
mechanism that operates on a node list \ and checks for an attribute can of
course be controlled by the macro package and some neat tricks in there, but just
avoiding a treatment when there is no demand also makes sense. This particular
optimization at first sight looks substantial but in the end only gives a few
percent gain, depending on the usage pattern. We just have to remember that over
the decades of \MKIV\ and \MKXL\ a lot already has been optimized.

Often playing with this itself leads to extra features, like tracing, so that we
get a bit more insight in what happens. In that case we accept the extra overhead
and accept a little less gain. Talking of tracing: one reason for the 2025
optimization was that the visualizer code has quite some impact on a run once
enabled. So, again here doesn't really benefit users, more ourselves. It'a also a
fact that there are many features that never kick in at the same time. If one
uses 20 attribute driven features, there is more gain to expect than when a 2 are
used. The reason is not so much in the tenfold usage but in the simple fact that
of the 20 most are very limited in scope. You don't want a single character
kerned word on the title page influence the whole document runtime. vertical
spacing on the other hand is always used. Most if the special math features are
seldom used so why check on them of no attribute is active? Those are the places
where we can gain a little, and all small amounts then add up. But, this is very
much a \CONTEXT\ thing.

% test attr list instead .. tricky housekeeping

\stopchapter

\stopcomponent
